\documentclass[aspectratio=169,10pt]{beamer}

% Shared Beamer setup for Fall 2026 course slides.
\mode<presentation>{
  \usetheme{Madrid}
  \usecolortheme{default}
}

\definecolor{CalPolyGreen}{HTML}{154734}
\definecolor{CalPolyGold}{HTML}{BD8B13}
\definecolor{DeckInk}{HTML}{1F2933}
\definecolor{DeckMuted}{HTML}{5F6B7A}

\setbeamercolor{structure}{fg=CalPolyGreen}
\setbeamercolor{palette primary}{bg=CalPolyGreen,fg=white}
\setbeamercolor{palette secondary}{bg=DeckInk,fg=white}
\setbeamercolor{palette tertiary}{bg=CalPolyGold,fg=DeckInk}
\setbeamercolor{frametitle}{bg=CalPolyGreen,fg=white}
\setbeamercolor{title}{fg=CalPolyGreen}
\setbeamercolor{normal text}{fg=DeckInk,bg=white}
\setbeamercolor{block title}{bg=CalPolyGreen,fg=white}
\setbeamercolor{block body}{bg=black!3,fg=DeckInk}

\setbeamertemplate{navigation symbols}{}
\setbeamertemplate{footline}[frame number]
\setbeamertemplate{itemize item}{\color{CalPolyGold}$\blacktriangleright$}
\setbeamertemplate{itemize subitem}{\color{CalPolyGreen}$\bullet$}

\newcommand{\CourseNumber}{Course}
\newcommand{\CourseTitle}{Course Title}
\newcommand{\LectureNumber}{0}
\newcommand{\LectureTitle}{Lecture Title}
\newcommand{\LectureDate}{Date}
\newcommand{\CourseUnit}{Unit}

\newcommand{\coursenumber}[1]{\renewcommand{\CourseNumber}{#1}}
\newcommand{\coursetitle}[1]{\renewcommand{\CourseTitle}{#1}}
\newcommand{\lecturenumber}[1]{\renewcommand{\LectureNumber}{#1}}
\newcommand{\lecturetitle}[1]{\renewcommand{\LectureTitle}{#1}}
\newcommand{\lecturedate}[1]{\renewcommand{\LectureDate}{#1}}
\newcommand{\courseunit}[1]{\renewcommand{\CourseUnit}{#1}}

\author{Kirk Duran}
\institute{California Polytechnic State University, San Luis Obispo}
\date{\LectureDate}

\newcommand{\makedecktitle}{
  \begin{frame}[plain]
    \vspace{0.25in}
    {\usebeamerfont{title}\color{CalPolyGreen}\LectureTitle\par}
    \vspace{0.18in}
    {\large \CourseNumber{}: \CourseTitle\par}
    \vspace{0.08in}
    {\large Lecture \LectureNumber{} -- \LectureDate\par}
    \vfill
    {\small Kirk Duran \textbar{} Cal Poly, San Luis Obispo\par}
  \end{frame}
}

\newcommand{\placeholdernote}{
  \begin{block}{Placeholder}
    Replace these bullets with the final lecture material, examples, and in-class activities.
  \end{block}
}

\AtBeginSection[]{
  \begin{frame}{Roadmap}
    \tableofcontents[currentsection]
  \end{frame}
}


\graphicspath{{../assets/lecture-05/}}

% If an image file is not yet available, the slide displays a labeled
% placeholder using the intended filename. Place the matching file in
% ../assets/lecture-05/ to replace the placeholder automatically.
\newcommand{\lectureimage}[2][1.75in]{%
  \IfFileExists{../assets/lecture-05/#2}{%
    \includegraphics[width=\linewidth,height=#1,keepaspectratio]{#2}%
  }{%
    \fbox{%
      \parbox[c][#1][c]{0.94\linewidth}{%
        \centering
        \small\textbf{Image Placeholder}\\[0.08in]
        \scriptsize\texttt{\detokenize{#2}}
      }%
    }%
  }%
}

\setbeamersize{text margin left=0.42in,text margin right=0.42in}

\coursenumber{CS 1001}
\coursetitle{Introduction to Computing}
\lecturenumber{5}
\lecturetitle{Variables, Assignment, and Program State}
\lecturedate{September 2, 2026}
\courseunit{1. Computing and Program Execution}

\begin{document}

\makedecktitle

\begin{frame}{Today}

  \begin{itemize}

    \item How a program gives a value a name and uses that value later.

    \item Why assignment evaluates the right-hand side before changing program state.

    \item How Python connects names, namespaces, objects, and object identities.

    \item How reassignment, dynamic typing, and immutable strings fit this model.

    \item How comments, annotations, and good variable names make state easier to understand.

  \end{itemize}

  \begin{keyidea}

    This is the first lecture in which the order of multiple statements changes what later statements mean.

  \end{keyidea}

\end{frame}

\sectionframe{Part 1: From Expressions to Program State}

\begin{frame}{Expressions Produce Values; Programs Can Retain Them}

  \begin{columns}[T,onlytextwidth]

    \begin{column}{0.52\textwidth}

      \begin{block}{Lecture 3}

        \texttt{10 + 5} evaluates to one value, \texttt{15}. Nothing in the expression gives that result a name for later use.

      \end{block}

      \begin{block}{Lecture 5}

        \texttt{score = 10 + 5} evaluates the same expression and then associates its result with the name \texttt{score}.

      \end{block}

    \end{column}

    \begin{column}{0.42\textwidth}

      % Search direction: isolated expression contrasted with a sequence that retains a value.
      \lectureimage[2.24in]{expression-versus-stateful-program.png}

    \end{column}

  \end{columns}

  \begin{keyidea}

    A named result can influence statements that execute later.

  \end{keyidea}

\end{frame}

\begin{frame}{A Variable Is a Name Associated with a Value}

  \begin{cpdefinition}

    A \textbf{variable} is a name that is currently associated with a value.

  \end{cpdefinition}

  \begin{columns}[T,onlytextwidth]

    \begin{column}{0.47\textwidth}

      \begin{block}{Statement}

        \centering
        \Large\texttt{score = 10}

      \end{block}

    \end{column}

    \begin{column}{0.47\textwidth}

      \begin{block}{Association after execution}

        \centering
        \Large\texttt{score} $\longrightarrow$ \texttt{10}

      \end{block}

    \end{column}

  \end{columns}

  \vspace{0.10in}

  \begin{itemize}

    \item \texttt{score} is the name; \texttt{10} is the value.

    \item Later, evaluating \texttt{score} retrieves its current value.

    \item The name and the value are related, but they are not the same thing.

  \end{itemize}

\end{frame}

\begin{frame}{Assignment Evaluates First and Associates Second}

  \begin{columns}[T,onlytextwidth]

    \begin{column}{0.54\textwidth}

      \begin{block}{Statement}

        \centering
        \Large\texttt{score = 7 + 3}

      \end{block}

      \begin{enumerate}

        \item Evaluate the right-hand expression: \texttt{7 + 3} $\longrightarrow$ \texttt{10}.

        \item Associate the name \texttt{score} with the resulting value \texttt{10}.

      \end{enumerate}

    \end{column}

    \begin{column}{0.40\textwidth}

      % Search direction: two-stage assignment flow, evaluate right side then bind left-side name.
      \lectureimage[2.30in]{assignment-evaluate-then-bind.png}

    \end{column}

  \end{columns}

  \begin{keyidea}

    The value must exist before Python can associate the name with it.

  \end{keyidea}

\end{frame}

\begin{frame}{Assignment Is Not Algebraic Equality}

  \begin{columns}[T,onlytextwidth]

    \begin{column}{0.48\textwidth}

      \begin{block}{In algebra}

        $x = x + 1$ would assert that two quantities are equal, which is impossible for ordinary numbers.

      \end{block}

    \end{column}

    \begin{column}{0.48\textwidth}

      \begin{block}{In Python}

        \texttt{x = x + 1} is an instruction:

        \begin{enumerate}
          \item retrieve the current value of \texttt{x};
          \item add \texttt{1};
          \item associate \texttt{x} with the new result.
        \end{enumerate}

      \end{block}

    \end{column}

  \end{columns}

  \begin{keyidea}

    Read \texttt{=} as ``assign'' or ``associate,'' not ``is mathematically equal to.''

  \end{keyidea}

\end{frame}

\begin{frame}{Program State Is a Snapshot of Current Associations}

  \begin{cpdefinition}

    A program's \textbf{state} is the collection of variable--value associations that currently exist.

  \end{cpdefinition}

  \begin{columns}[T,onlytextwidth]

    \begin{column}{0.44\textwidth}

      \begin{block}{Statements executed}

        \texttt{score = 10}\\
        \texttt{bonus = 5}

      \end{block}

    \end{column}

    \begin{column}{0.50\textwidth}

      \begin{center}
      \begin{tabular}{l|c}
        \textbf{Name} & \textbf{Current value} \\\hline
        \texttt{score} & \texttt{10} \\
        \texttt{bonus} & \texttt{5}
      \end{tabular}
      \end{center}

    \end{column}

  \end{columns}

  \begin{keyidea}

    Executing an assignment can change the program's state; merely evaluating an expression does not create a new name.

  \end{keyidea}

\end{frame}

\begin{frame}{Reassignment Replaces a Name's Current Association}

  \begin{columns}[T,onlytextwidth]

    \begin{column}{0.51\textwidth}

      \begin{block}{Sequence}

        \texttt{score = 10}\\
        \texttt{score = 25}

      \end{block}

      \begin{itemize}

        \item After line 1: \texttt{score} $\longrightarrow$ \texttt{10}

        \item After line 2: \texttt{score} $\longrightarrow$ \texttt{25}

        \item The name has one current association in this simple model.

      \end{itemize}

    \end{column}

    \begin{column}{0.43\textwidth}

      % Search direction: label score moving from object 10 to object 25.
      \lectureimage[2.23in]{reassignment-moves-binding.png}

    \end{column}

  \end{columns}

\end{frame}

\sectionframe{Part 2: Names, Objects, and Memory}

\begin{frame}{Python Connects Names to Objects}

  \begin{columns}[T,onlytextwidth]

    \begin{column}{0.54\textwidth}

      \begin{itemize}

        \item An assignment \textbf{binds a name} to an object.

        \item The name is recorded in a \textbf{namespace}.

        \item The object carries a value, a type, and an identity.

        \item Evaluating the name asks its namespace which object is currently associated with that name.

      \end{itemize}

      \begin{cpdefinition}

        A \textbf{namespace} is a mapping from names to objects.

      \end{cpdefinition}

    \end{column}

    \begin{column}{0.40\textwidth}

      % Search direction: namespace table mapping labels score and message to distinct Python objects.
      \lectureimage[2.28in]{namespace-name-object-mapping.png}

    \end{column}

  \end{columns}

  \note{[Sources] Python Language Reference, ``Naming and binding,'' https://docs.python.org/3/reference/executionmodel.html\#naming-and-binding.}

\end{frame}

\begin{frame}{Where Is the Variable Name Stored?}

  \begin{columns}[T,onlytextwidth]

    \begin{column}{0.53\textwidth}

      \begin{itemize}

        \item At the top level of a Colab cell or Python file, assignment records a name in the module's global namespace.

        \item Python exposes a module namespace as a dictionary-like mapping.

        \item Conceptually, an entry looks like:

      \end{itemize}

      \begin{block}{Namespace entry}

        \centering
        \texttt{"score"} $\longrightarrow$ object representing \texttt{10}

      \end{block}

      \begin{exampleblockcp}

        The spelling \texttt{score} is stored separately from the numerical information in the integer object.

      \end{exampleblockcp}

    \end{column}

    \begin{column}{0.41\textwidth}

      % Search direction: variable name text stored as a namespace key pointing to a separate object.
      \lectureimage[2.30in]{variable-name-storage-namespace.png}

    \end{column}

  \end{columns}

  \note{[Sources] Python Language Reference, ``Modules,'' https://docs.python.org/3/reference/datamodel.html\#modules.}

\end{frame}

\begin{frame}{Every Python Object Has Identity, Type, and Value}

  \begin{columns}[T,onlytextwidth]

    \begin{column}{0.55\textwidth}

      \begin{description}

        \item[Identity] Distinguishes this object from every other object during its lifetime.

        \item[Type] Determines which operations the object supports and helps determine how its value is represented.

        \item[Value] The information represented by the object, such as \texttt{10} or \texttt{"hello"}.

      \end{description}

      \begin{keyidea}

        The type belongs to the object. A name can later be associated with an object of another type.

      \end{keyidea}

    \end{column}

    \begin{column}{0.39\textwidth}

      % Search direction: one Python object card labeled identity, type, and value.
      \lectureimage[2.24in]{object-identity-type-value.png}

    \end{column}

  \end{columns}

  \note{[Sources] Python Language Reference, ``Objects, values and types,'' https://docs.python.org/3/reference/datamodel.html\#objects-values-and-types.}

\end{frame}

\begin{frame}{\texttt{id()} Observes the Current Object's Identity}

  \begin{columns}[T,onlytextwidth]

    \begin{column}{0.55\textwidth}

      \begin{block}{Colab experiment}

        \texttt{score = 10}\\
        \texttt{id(score)}

      \end{block}

      \begin{itemize}

        \item \texttt{id(score)} first evaluates \texttt{score} to retrieve its object.

        \item It then returns an integer that identifies that object during its lifetime.

        \item The particular integer is not meaningful to us and can differ across runs.

        \item It is the identity of the \textbf{object}, not an ``address of the variable name.''

      \end{itemize}

    \end{column}

    \begin{column}{0.39\textwidth}

      % Search direction: Colab id(score) output pointing to object identity rather than variable label.
      \lectureimage[2.30in]{id-current-object-cpython-address.png}

    \end{column}

  \end{columns}

  \note{[Sources] Python built-in function id(), https://docs.python.org/3/library/functions.html\#id; Python Data Model, https://docs.python.org/3/reference/datamodel.html\#objects-values-and-types.}

\end{frame}

\begin{frame}{Is an Object's Identity Its Memory Address?}

  \begin{columns}[T,onlytextwidth]

    \begin{column}{0.47\textwidth}

      \begin{block}{Python's language-level promise}

        \texttt{id()} returns an integer that is unique and constant for the object during its lifetime.

      \end{block}

    \end{column}

    \begin{column}{0.47\textwidth}

      \begin{block}{CPython's implementation}

        In CPython, that integer is the object's memory address. Other Python implementations need not use an address.

      \end{block}

    \end{column}

  \end{columns}

  \begin{itemize}

    \item Memory management is automatic; programmers ordinarily do not choose these addresses.

    \item After an object no longer exists, a later object may reuse an identity value.

    \item Use \texttt{id()} to test the mental model, not to build arithmetic or program logic around the number.

  \end{itemize}

  \begin{keyidea}

    ``Identity'' is the portable idea; ``memory address'' is a CPython-specific implementation detail.

  \end{keyidea}

  \note{[Sources] Python built-in function id(), https://docs.python.org/3/library/functions.html\#id.}

\end{frame}

\begin{frame}{Assignment Changes the Namespace Entry}

  \begin{columns}[T,onlytextwidth]

    \begin{column}{0.47\textwidth}

      \begin{block}{Before reassignment}

        \texttt{score = 10}\\[0.08in]
        namespace: \texttt{score} $\longrightarrow$ integer object \texttt{10}

      \end{block}

    \end{column}

    \begin{column}{0.47\textwidth}

      \begin{block}{After reassignment}

        \texttt{score = 25}\\[0.08in]
        namespace: \texttt{score} $\longrightarrow$ integer object \texttt{25}

      \end{block}

    \end{column}

  \end{columns}

  \vspace{0.10in}

  \begin{keyidea}

    Reassignment does not edit the integer \texttt{10}. It changes which object the name \texttt{score} denotes.

  \end{keyidea}

  \begin{activity}

    Predict whether \texttt{id(score)} must remain the same after \texttt{score = 25}; then test the prediction in Colab.

  \end{activity}

\end{frame}

\begin{frame}{String Reassignment Creates a New Result}

  \begin{columns}[T,onlytextwidth]

    \begin{column}{0.52\textwidth}

      \begin{block}{Sequence}

        \texttt{message = "hello"}\\
        \texttt{message = message + " world"}

      \end{block}

      \begin{enumerate}

        \item Retrieve the current string \texttt{"hello"}.

        \item Concatenate it with \texttt{" world"}.

        \item Produce a new string \texttt{"hello world"}.

        \item Rebind \texttt{message} to the new result.

      \end{enumerate}

    \end{column}

    \begin{column}{0.42\textwidth}

      % Search direction: message label moving from immutable hello string to a new hello world string.
      \lectureimage[2.38in]{string-reassignment-new-object.png}

    \end{column}

  \end{columns}

\end{frame}

\begin{frame}{Why Are Strings Immutable?}

  \begin{cpdefinition}

    A string is \textbf{immutable}: after a string object is created, its textual value cannot be altered.

  \end{cpdefinition}

  \begin{columns}[T,onlytextwidth]

    \begin{column}{0.52\textwidth}

      \begin{itemize}

        \item Text often represents identifiers, usernames, passwords, filenames, and dictionary or database keys.

        \item Stable string values are easier to share and reason about safely.

        \item Operations that appear to modify text instead produce another string value.

      \end{itemize}

      \begin{exampleblockcp}

        \texttt{message += " world"} still produces a new string and rebinds \texttt{message} for the simple string case.

      \end{exampleblockcp}

    \end{column}

    \begin{column}{0.42\textwidth}

      % Search direction: immutable string used as stable identifier and database key.
      \lectureimage[2.30in]{string-immutability-stable-identifiers.png}

    \end{column}

  \end{columns}

  \note{[Sources] Python Language Reference, immutable sequence types, https://docs.python.org/3/library/stdtypes.html\#immutable-sequence-types.}

\end{frame}

\sectionframe{Part 3: Special Values and Dynamic Typing}

\begin{frame}{\texttt{None} Represents the Intentional Absence of a Value}

  \begin{columns}[T,onlytextwidth]

    \begin{column}{0.54\textwidth}

      \begin{block}{Example}

        \texttt{result = None}

      \end{block}

      \begin{itemize}

        \item \texttt{None} is a literal value, not an unassigned name.

        \item Its type is \texttt{NoneType}.

        \item Python provides exactly one \texttt{None} object, so it is a singleton.

        \item A variable associated with \texttt{None} \emph{has} a value: the value representing absence.

      \end{itemize}

      \begin{exampleblockcp}

        \texttt{type(None)} $\longrightarrow$ \texttt{NoneType}

      \end{exampleblockcp}

    \end{column}

    \begin{column}{0.40\textwidth}

      % Search direction: many variable labels pointing to the single shared Python None object.
      \lectureimage[2.30in]{none-singleton-object.png}

    \end{column}

  \end{columns}

  \note{[Sources] Python Language Reference, ``None,'' https://docs.python.org/3/reference/datamodel.html\#none.}

\end{frame}

\begin{frame}{The ``Billion-Dollar Mistake'' Is a Warning About Absence}

  \begin{columns}[T,onlytextwidth]

    \begin{column}{0.56\textwidth}

      \begin{itemize}

        \item Computer scientist Tony Hoare introduced a null reference into ALGOL W's type system in 1965.

        \item In a 2009 talk, he called that decision his ``billion-dollar mistake,'' citing errors, vulnerabilities, and crashes caused by unchecked null references.

        \item Python's \texttt{None} is not identical to every language's null mechanism, but it carries the same general lesson: absence must be represented deliberately and handled explicitly.

      \end{itemize}

      \begin{keyidea}

        A value meaning ``nothing here'' is still part of program state and must be reasoned about.

      \end{keyidea}

    \end{column}

    \begin{column}{0.38\textwidth}

      % Search direction: portrait of Tony Hoare with ALGOL W and null-reference historical context.
      \lectureimage[2.35in]{tony-hoare-null-reference.png}

    \end{column}

  \end{columns}

  \note{[Sources] Tony Hoare, ``Null References: The Billion Dollar Mistake,'' QCon London 2009, https://qconlondon.com/london-2009/qconlondon.com/london-2009/presentation/Null\%2BReferences\_\%2BThe\%2BBillion\%2BDollar\%2BMistake.html.}

\end{frame}

\begin{frame}{Python Is Dynamically Typed}

  \begin{columns}[T,onlytextwidth]

    \begin{column}{0.50\textwidth}

      \begin{block}{The same name can be rebound}

        \texttt{item = 42}\hfill integer object\\
        \texttt{item = "hello"}\hfill string object\\
        \texttt{item = 4.0}\hfill floating-point object

      \end{block}

      \begin{itemize}

        \item Each object has a definite type.

        \item The name \texttt{item} does not have one permanent type.

        \item Its current value and type depend on the most recent executed assignment.

      \end{itemize}

    \end{column}

    \begin{column}{0.44\textwidth}

      % Search direction: one variable name successively pointing to int, str, and float objects.
      \lectureimage[2.35in]{dynamic-rebinding-object-types.png}

    \end{column}

  \end{columns}

\end{frame}

\begin{frame}{Dynamic Typing Has Benefits and Costs}

  \begin{columns}[T,onlytextwidth]

    \begin{column}{0.47\textwidth}

      \begin{block}{Benefits}

        \begin{itemize}
          \item less type declaration syntax;
          \item rapid experimentation;
          \item flexible reuse of code.
        \end{itemize}

      \end{block}

    \end{column}

    \begin{column}{0.47\textwidth}

      \begin{block}{Costs}

        \begin{itemize}
          \item some type mismatches appear only when code runs;
          \item a name's intended kind of value may be less obvious;
          \item changes in state require careful reading.
        \end{itemize}

      \end{block}

    \end{column}

  \end{columns}

  \begin{exampleblockcp}

    \texttt{item = 42} followed by \texttt{item = "hello"} is legal, but a reader must understand whether that change is intentional.

  \end{exampleblockcp}

\end{frame}

\begin{frame}{Type Annotations Communicate Intended Types}

  \begin{columns}[T,onlytextwidth]

    \begin{column}{0.53\textwidth}

      \begin{block}{Annotated assignment}

        \texttt{score: int = 42}

      \end{block}

      \begin{itemize}

        \item \texttt{: int} records that \texttt{score} is intended to refer to an integer.

        \item Editors, documentation tools, and type checkers can use the annotation.

        \item Standard Python execution ordinarily does not enforce the annotation by itself.

      \end{itemize}

      \begin{exampleblockcp}

        An annotation supports communication and analysis; assignment still determines the current runtime association.

      \end{exampleblockcp}

    \end{column}

    \begin{column}{0.41\textwidth}

      % Search direction: annotated Python variable with editor or type-checker feedback.
      \lectureimage[2.25in]{type-annotation-intent-not-enforcement.png}

    \end{column}

  \end{columns}

  \note{[Sources] Python documentation, typing specification and runtime behavior, https://docs.python.org/3/library/typing.html.}

\end{frame}

\sectionframe{Part 4: Updating State}

\begin{frame}{Updating a Variable Uses Its Current Value}

  \begin{block}{Statement}

    \centering
    \Large\texttt{score = score + 5}

  \end{block}

  \begin{enumerate}

    \item Look up the current object associated with \texttt{score}.

    \item Evaluate the right-hand expression \texttt{score + 5}.

    \item Associate \texttt{score} with the resulting object.

  \end{enumerate}

  \begin{keyidea}

    Every use of \texttt{score} on the right occurs before the assignment changes \texttt{score} on the left.

  \end{keyidea}

\end{frame}

\begin{frame}{A Variable Retains a Result, Not the Formula}

  \begin{columns}[T,onlytextwidth]

    \begin{column}{0.49\textwidth}

      \begin{block}{Sequence}

        \texttt{x = 4}\\
        \texttt{y = x + 2}\\
        \texttt{x = 10}

      \end{block}

      \begin{activity}

        What are the final values of \texttt{x} and \texttt{y}? Does the final assignment to \texttt{x} retroactively change \texttt{y}?

      \end{activity}

    \end{column}

    \begin{column}{0.45\textwidth}

      % Search direction: x and y object-binding diagram showing y remains 6 after x changes to 10.
      \lectureimage[2.38in]{xy-reassignment-no-retroactive-change.png}

    \end{column}

  \end{columns}

  \begin{keyidea}

    \texttt{y = x + 2} evaluates once. The name \texttt{y} is associated with the resulting value \texttt{6}, not with a live formula.

  \end{keyidea}

\end{frame}

\begin{frame}{Accumulation Repeatedly Updates One Association}

  \begin{columns}[T,onlytextwidth]

    \begin{column}{0.49\textwidth}

      \begin{block}{Sequence}

        \texttt{total = 0}\\
        \texttt{total = total + 5}\\
        \texttt{total = total + 8}

      \end{block}

    \end{column}

    \begin{column}{0.45\textwidth}

      \begin{center}
      \renewcommand{\arraystretch}{1.35}
      \begin{tabular}{l|c}
        \textbf{After line} & \textbf{\texttt{total}} \\\hline
        1 & \texttt{0} \\
        2 & \texttt{5} \\
        3 & \texttt{13}
      \end{tabular}
      \end{center}

    \end{column}

  \end{columns}

  \begin{keyidea}

    Each update uses the current state to calculate the next state.

  \end{keyidea}

\end{frame}

\begin{frame}{A Name Must Be Assigned Before It Is Evaluated}

  \begin{block}{No current association}

    \texttt{points + 5}

  \end{block}

  \begin{itemize}

    \item Python must retrieve \texttt{points} before it can perform addition.

    \item If no accessible binding for \texttt{points} exists, evaluation cannot continue.

    \item Python reports a \texttt{NameError}.

  \end{itemize}

  \begin{exampleblockcp}

    Correct the state first: execute \texttt{points = 0}, then evaluate or update \texttt{points}.

  \end{exampleblockcp}

  \note{[Sources] Python Language Reference, ``Resolution of names,'' https://docs.python.org/3/reference/executionmodel.html\#resolution-of-names.}

\end{frame}

\begin{frame}{Augmented Assignment Is a Compact Update}

  \begin{columns}[T,onlytextwidth]

    \begin{column}{0.50\textwidth}

      \begin{block}{Ordinary reassignment}

        \texttt{total = total + 5}

      \end{block}

      \begin{block}{Augmented assignment}

        \texttt{total += 5}

      \end{block}

      \begin{itemize}

        \item Retrieve the current \texttt{total}.

        \item Add \texttt{5}.

        \item Associate \texttt{total} with the result.

      \end{itemize}

    \end{column}

    \begin{column}{0.44\textwidth}

      % Search direction: ordinary reassignment and augmented assignment shown as the same numeric update flow.
      \lectureimage[2.35in]{augmented-assignment-update.png}

    \end{column}

  \end{columns}

  \begin{exampleblockcp}

    For today's integer, float, and string examples, this is the appropriate mental shorthand. Later data structures require additional care.

  \end{exampleblockcp}

\end{frame}

\begin{frame}{Common Augmented Assignment Operators}

  \begin{center}
  \renewcommand{\arraystretch}{1.25}
  \begin{tabular}{c|l@{\qquad}c|l}
    \textbf{Operator} & \textbf{Update} & \textbf{Operator} & \textbf{Update} \\\hline
    \texttt{+=} & add, then reassign & \texttt{-=} & subtract, then reassign \\
    \texttt{*=} & multiply, then reassign & \texttt{/=} & divide, then reassign \\
    \texttt{//=} & floor-divide, then reassign & \texttt{\%=} & take modulus, then reassign \\
    \texttt{**=} & exponentiate, then reassign & &
  \end{tabular}
  \end{center}

  \vspace{0.14in}

  \begin{exampleblockcp}

    If \texttt{count} currently refers to \texttt{8}, then \texttt{count *= 3} associates \texttt{count} with \texttt{24}.

  \end{exampleblockcp}

  \begin{keyidea}

    Introduce the compact form only after you can explain the ordinary reassignment it represents.

  \end{keyidea}

\end{frame}

\sectionframe{Part 5: Readable State and Deliberate Practice}

\begin{frame}{Comments Explain Code Without Changing State}

  \begin{columns}[T,onlytextwidth]

    \begin{column}{0.54\textwidth}

      \begin{block}{End-of-line comments}

        \texttt{score = 10 \# starting score}\\
        \texttt{score += 5 \# add the bonus}

      \end{block}

      \begin{itemize}

        \item A comment begins with \texttt{\#} and continues to the end of that line.

        \item Python ignores the comment when executing the statement.

        \item A useful comment explains purpose or reasoning, not merely the visible syntax.

      \end{itemize}

    \end{column}

    \begin{column}{0.40\textwidth}

      % Search direction: Python code with comment text visually separated from state-changing statements.
      \lectureimage[2.28in]{comments-versus-state-changes.png}

    \end{column}

  \end{columns}

  \begin{alertblock}{Technical distinction}

    Triple-quoted text creates a string literal. It is sometimes used for documentation, but it is not Python's comment syntax.

  \end{alertblock}

  \note{[Sources] Python Language Reference, ``Comments,'' https://docs.python.org/3/reference/lexical_analysis.html\#comments.}

\end{frame}

\begin{frame}[shrink=8]{Variable Naming: Valid, Descriptive, and Conventional}

  \begin{columns}[T,onlytextwidth]

    \begin{column}{0.52\textwidth}

      \begin{block}{Python's basic identifier rules}

        \begin{itemize}
          \item use letters, digits, and underscores;
          \item do not begin with a digit;
          \item do not include spaces or hyphens;
          \item capitalization matters;
          \item do not use a reserved keyword.
        \end{itemize}

      \end{block}

      \begin{block}{Course convention}

        Use lowercase words separated by underscores: \texttt{course\_units}, \texttt{total\_cost}, \texttt{quiz\_score}.

      \end{block}

    \end{column}

    \begin{column}{0.42\textwidth}

      \begin{block}{Reserved keywords}

        Python reserves words such as \texttt{if}, \texttt{else}, \texttt{for}, \texttt{while}, \texttt{class}, \texttt{def}, \texttt{return}, \texttt{import}, \texttt{True}, \texttt{False}, and \texttt{None}. They cannot be variable names.

      \end{block}

      \begin{exampleblockcp}

        Prefer \texttt{exam\_score = 92}, not \texttt{x = 92}, when the value's role matters.

      \end{exampleblockcp}

    \end{column}

  \end{columns}

  \note{[Sources] Python Language Reference, ``Identifiers and keywords,'' https://docs.python.org/3/reference/lexical_analysis.html\#identifiers; PEP 8, naming conventions, https://peps.python.org/pep-0008/\#naming-conventions.}

\end{frame}

\begin{frame}{Keywords and Built-in Names Are Different}

  \begin{columns}[T,onlytextwidth]

    \begin{column}{0.47\textwidth}

      \begin{block}{Reserved keywords}

        \texttt{if = 3} is invalid syntax because \texttt{if} is part of Python's grammar.

      \end{block}

      \begin{itemize}

        \item Keywords cannot be used as variable names.

        \item Python publishes the current keyword list through the \texttt{keyword} module.

      \end{itemize}

    \end{column}

    \begin{column}{0.47\textwidth}

      \begin{block}{Built-in names}

        \texttt{print = 3} is legal, but it hides the built-in name \texttt{print} in the current namespace.

      \end{block}

      \begin{itemize}

        \item Avoid names such as \texttt{print}, \texttt{type}, \texttt{id}, \texttt{int}, and \texttt{str}.

        \item Legal does not always mean readable or safe.

      \end{itemize}

    \end{column}

  \end{columns}

  \begin{keyidea}

    A keyword is forbidden by the grammar; a built-in name is merely easy to shadow accidentally.

  \end{keyidea}

\end{frame}

\begin{frame}{Colab Lab: Observe Rebinding with \texttt{id()}}

  \begin{columns}[T,onlytextwidth]

    \begin{column}{0.52\textwidth}

      \begin{block}{Run one line at a time}

        \texttt{message = "hello"}\\
        \texttt{id(message)}\\
        \texttt{message += " world"}\\
        \texttt{id(message)}

      \end{block}

      \begin{activity}

        Before running the second \texttt{id()} call, predict whether the object identity will be the same. Explain the observation using ``evaluate,'' ``new string,'' ``namespace,'' and ``rebind.''

      \end{activity}

    \end{column}

    \begin{column}{0.42\textwidth}

      % Search direction: Google Colab cells comparing id(message) before and after string augmentation.
      \lectureimage[2.42in]{colab-id-string-reassignment.png}

    \end{column}

  \end{columns}

  \begin{alertblock}{Interpretation rule}

    Compare whether the identity changed; do not try to interpret the magnitude of the integer.

  \end{alertblock}

\end{frame}

\begin{frame}{Misconceptions to Correct Explicitly}

  \begin{itemize}

    \item \texttt{x = x + 1} is valid because assignment is an action, not an algebraic claim.

    \item The entire right-hand side is evaluated before the left-hand name is updated.

    \item A variable retains the resulting object, not the expression that produced it.

    \item Reassigning \texttt{x} does not retroactively update \texttt{y}.

    \item A name must be assigned before Python can evaluate it.

    \item \texttt{id(x)} identifies the current object, not the storage location of the written name \texttt{x}.

    \item A name can be rebound to an object of another type; the objects themselves still have definite types.

  \end{itemize}

\end{frame}

\begin{frame}{Exit Ticket: Explain the Final State}

  \begin{columns}[T,onlytextwidth]

    \begin{column}{0.43\textwidth}

      \begin{block}{Statements}

        \texttt{level = 2}\\
        \texttt{next\_level = level + 1}\\
        \texttt{level = 10}\\
        \texttt{level += 2}

      \end{block}

    \end{column}

    \begin{column}{0.51\textwidth}

      \begin{activity}

        Report the final values of \texttt{level} and \texttt{next\_level}. Then explain:

        \begin{enumerate}
          \item what each assignment evaluates;
          \item which namespace association changes;
          \item why \texttt{next\_level} does not change later.
        \end{enumerate}

      \end{activity}

    \end{column}

  \end{columns}

  \begin{keyidea}

    Outcome: you can predict how a sequence of assignments changes program state and explain that state using names, objects, and values.

  \end{keyidea}

\end{frame}

\begin{frame}{Image Placeholder Index}

  \scriptsize
  \begin{columns}[T,onlytextwidth]

    \begin{column}{0.49\textwidth}

      \begin{itemize}
        \item \texttt{expression-versus-stateful-program.png}
        \item \texttt{assignment-evaluate-then-bind.png}
        \item \texttt{reassignment-moves-binding.png}
        \item \texttt{namespace-name-object-mapping.png}
        \item \texttt{variable-name-storage-namespace.png}
        \item \texttt{object-identity-type-value.png}
        \item \texttt{id-current-object-cpython-address.png}
        \item \texttt{string-reassignment-new-object.png}
        \item \texttt{string-immutability-stable-identifiers.png}
      \end{itemize}

    \end{column}

    \begin{column}{0.49\textwidth}

      \begin{itemize}
        \item \texttt{none-singleton-object.png}
        \item \texttt{tony-hoare-null-reference.png}
        \item \texttt{dynamic-rebinding-object-types.png}
        \item \texttt{type-annotation-intent-not-enforcement.png}
        \item \texttt{xy-reassignment-no-retroactive-change.png}
        \item \texttt{augmented-assignment-update.png}
        \item \texttt{comments-versus-state-changes.png}
        \item \texttt{colab-id-string-reassignment.png}
      \end{itemize}

    \end{column}

  \end{columns}

  \begin{exampleblockcp}

    Save each image under \texttt{../assets/lecture-05/}. The labeled placeholder will be replaced automatically when its filename matches.

  \end{exampleblockcp}

\end{frame}

\end{document}
